```latex
\documentclass[12pt]{article}
\input{article_style.tex}

\title{Raport projektowy\\Uczenie się ze wzmocnieniem}
\author{Sebastian Gęborys}
\date{2026/2027}

\begin{document}
\maketitle
\tableofcontents
\newpage

\section{Informacje podstawowe}

\begin{center}
\begin{tabular}{p{0.34\textwidth}p{0.56\textwidth}}
\toprule
Element & Informacja \\
\midrule
Imię i nazwisko & Sebastian Gęborys \\
Numer albumu & 179487 \\
Kierunek/specjalność & Sztuczna inteligencja \\
Tytuł projektu & Implementacja i analiza agenta wykorzystującego uczenie ze wzmocnieniem do autonomicznego sterowania pojazdem w grze TrackMania 2020 z wykorzystaniem biblioteki TrackmaniaRL \\
Link do repozytorium & Do uzupełnienia \\
Data oddania & Do uzupełnienia \\
\bottomrule
\end{tabular}
\end{center}

\section{Temat i cel projektu}

\subsection{Opis problemu}

Przedmiotem projektu jest opracowanie agenta wykorzystującego uczenie ze wzmocnieniem (Reinforcement Learning) do autonomicznego sterowania pojazdem w grze TrackMania 2020.

Zadaniem agenta będzie samodzielne prowadzenie samochodu po wyznaczonej trasie poprzez podejmowanie decyzji dotyczących przyspieszania, hamowania oraz skręcania. Agent będzie otrzymywał informacje telemetryczne opisujące aktualną sytuację pojazdu i jego położenie względem trasy. Na podstawie tych informacji będzie wybierał odpowiednie działania sterujące.

Do komunikacji z grą oraz pozyskiwania danych telemetrycznych zostanie wykorzystana biblioteka TrackmaniaRL.

Problem polega na opracowaniu strategii jazdy pozwalającej agentowi na skuteczne pokonywanie trasy, ograniczenie błędów oraz uzyskiwanie możliwie krótkiego czasu przejazdu.

\subsection{Cel agenta}

Głównym celem agenta jest nauczenie się autonomicznego sterowania pojazdem w grze TrackMania 2020 w taki sposób, aby możliwie szybko i skutecznie pokonywał wyznaczoną trasę.

Agent powinien nauczyć się:
\begin{itemize}
    \item dobierania odpowiedniego przyspieszenia i hamowania,
    \item wykonywania skrętów z uwzględnieniem geometrii trasy,
    \item dostosowywania prędkości do aktualnych warunków jazdy,
    \item unikania wypadania poza trasę i innych błędów jazdy,
    \item planowania kolejnych działań z uwzględnieniem ich przyszłych konsekwencji.
\end{itemize}

Działanie agenta będzie oceniane na podstawie czasu przejazdu, skuteczności ukończenia trasy, liczby błędów oraz postępu uzyskiwanego podczas kolejnych epizodów.

\subsection{Uzasadnienie zastosowania uczenia się ze wzmocnieniem}

Problem autonomicznego sterowania pojazdem w grze TrackMania 2020 można rozpatrywać jako zadanie uczenia się ze wzmocnieniem, ponieważ wymaga podejmowania sekwencyjnych decyzji wpływających na dalszy przebieg jazdy.

Agent w każdym kroku otrzymuje obserwację opisującą aktualną sytuację pojazdu, a następnie wybiera akcję sterującą. Wykonanie tej akcji zmienia położenie, prędkość oraz orientację samochodu, co wpływa na kolejne dostępne możliwości działania.

Decyzje podejmowane przez agenta mają konsekwencje wykraczające poza pojedynczy krok. Przykładowo zbyt duża prędkość przed zakrętem może utrudnić jego pokonanie, doprowadzić do wypadnięcia z trasy i uniemożliwić uzyskanie dobrego czasu przejazdu. Z kolei wcześniejsze ograniczenie prędkości może pozwolić na bezpieczniejsze pokonanie zakrętu i uzyskanie lepszego wyniku całego przejazdu.

Istotny jest zatem wynik długoterminowy, a nie tylko natychmiastowa korzyść wynikająca z pojedynczej akcji.

Zastosowanie uczenia ze wzmocnieniem pozwoli agentowi na samodzielne doskonalenie strategii jazdy poprzez wielokrotne interakcje ze środowiskiem i analizowanie otrzymywanych nagród.

Problem nie jest zwykłą klasyfikacją ani jednorazową optymalizacją, ponieważ agent musi stale reagować na zmieniający się stan pojazdu i podejmować kolejne decyzje w trakcie całego przejazdu.

\section{Formalizacja problemu RL/MDP}

Problem zostanie wstępnie opisany jako proces decyzyjny Markowa (MDP), w którym agent podejmuje działania w środowisku gry, otrzymuje obserwacje i nagrody, a następnie na ich podstawie uczy się skuteczniejszego sterowania pojazdem.

\subsection{Agent i środowisko}

Agentem jest algorytm uczenia ze wzmocnieniem odpowiedzialny za autonomiczne sterowanie pojazdem w grze TrackMania 2020.

Środowiskiem jest gra TrackMania 2020, w której pojazd porusza się po wyznaczonej trasie zgodnie z zaimplementowanym modelem fizyki gry.

Do komunikacji ze środowiskiem, odczytu danych telemetrycznych oraz przekazywania poleceń sterujących zostanie wykorzystana biblioteka TrackmaniaRL.

Interakcja agenta ze środowiskiem będzie przebiegać w następujący sposób:
\begin{enumerate}
    \item Środowisko udostępnia agentowi aktualną obserwację telemetryczną.
    \item Agent analizuje otrzymane informacje.
    \item Agent wybiera akcję sterującą.
    \item Wybrana akcja jest przekazywana do gry.
    \item Pojazd wykonuje odpowiedni ruch.
    \item Środowisko zwraca nową obserwację oraz wartość nagrody.
\end{enumerate}

Proces jest powtarzany do momentu zakończenia epizodu.

\subsection{Stan lub obserwacja}

Obserwacja agenta będzie opierać się na danych telemetrycznych udostępnianych przez bibliotekę TrackmaniaRL.

Wstępnie zakłada się wykorzystanie następujących informacji:

\begin{center}
\begin{tabular}{p{0.30\textwidth}p{0.60\textwidth}}
\toprule
Składowa stanu/obserwacji & Znaczenie \\
\midrule
Pozycja pojazdu & Aktualne położenie samochodu w przestrzeni gry. \\
Prędkość & Aktualna prędkość poruszania się pojazdu. \\
Orientacja pojazdu & Kierunek, w którym skierowany jest samochód. \\
Położenie względem trasy & Informacja o położeniu pojazdu względem przebiegu toru. \\
Geometria trasy & Informacje o najbliższym fragmencie trasy, w szczególności o zakrętach. \\
Postęp na trasie & Informacja o stopniu pokonania wyznaczonej trasy. \\
\bottomrule
\end{tabular}
\end{center}

Powyższy zestaw stanowi wstępną propozycję obserwacji. Ostateczny zakres zmiennych zostanie określony po sprawdzeniu rzeczywistych danych udostępnianych przez bibliotekę TrackmaniaRL.

\subsection{Akcje}

Agent będzie podejmował decyzje dotyczące sterowania pojazdem. Wstępnie przewiduje się wykorzystanie następujących działań:

\begin{center}
\begin{tabular}{p{0.30\textwidth}p{0.60\textwidth}}
\toprule
Akcja & Opis \\
\midrule
Przyspieszanie & Zwiększanie prędkości pojazdu poprzez użycie gazu. \\
Hamowanie & Zmniejszanie prędkości pojazdu poprzez użycie hamulca. \\
Skręt w lewo & Zmiana kierunku jazdy w lewą stronę. \\
Skręt w prawo & Zmiana kierunku jazdy w prawą stronę. \\
Brak skrętu & Utrzymywanie aktualnego kierunku jazdy. \\
\bottomrule
\end{tabular}
\end{center}

Dokładna reprezentacja przestrzeni akcji, w szczególności wybór pomiędzy akcjami dyskretnymi i ciągłymi, zostanie ustalona na etapie implementacji po weryfikacji możliwości biblioteki TrackmaniaRL.

\subsection{Funkcja nagrody}

Funkcja nagrody będzie określać jakość decyzji podejmowanych przez agenta. Jej głównym zadaniem będzie zachęcanie do szybkiego i prawidłowego pokonywania trasy oraz ograniczania błędów jazdy.

Wstępnie rozważane składniki funkcji nagrody obejmują:
\begin{itemize}
    \item nagrodę za postęp wzdłuż trasy,
    \item uwzględnienie czasu przejazdu,
    \item karę za wypadnięcie poza trasę,
    \item karę za kolizje i inne błędy jazdy,
    \item dodatkową nagrodę za ukończenie trasy.
\end{itemize}

Wstępna postać funkcji nagrody może zostać zapisana jako:

\begin{equation}
    r_t =
    w_p \Delta p_t
    - w_c C_t
    - w_o O_t
    - w_t \Delta t
    + w_f F_t
\end{equation}

gdzie:
\begin{itemize}
    \item $\Delta p_t$ oznacza postęp pojazdu na trasie,
    \item $C_t$ oznacza wystąpienie kolizji,
    \item $O_t$ oznacza wypadnięcie poza trasę,
    \item $\Delta t$ oznacza czas trwania kroku,
    \item $F_t$ oznacza ukończenie trasy,
    \item $w_p$, $w_c$, $w_o$, $w_t$, $w_f$ oznaczają współczynniki wagowe poszczególnych składników.
\end{itemize}

Przedstawiona funkcja ma charakter wstępny. Konkretne wartości współczynników oraz sposób obliczania poszczególnych składników zostaną ustalone i zweryfikowane podczas pierwszych testów środowiska.

\subsection{Warunki zakończenia epizodu}

Epizod będzie odpowiadał pojedynczej próbie pokonania wyznaczonej trasy.

Wstępnie przewiduje się zakończenie epizodu w następujących sytuacjach:
\begin{itemize}
    \item ukończenie trasy przez pojazd,
    \item wypadnięcie poza trasę uniemożliwiające kontynuowanie przejazdu,
    \item wystąpienie sytuacji uznanej za poważny błąd jazdy,
    \item przekroczenie ustalonego limitu czasu,
    \item osiągnięcie maksymalnej liczby kroków epizodu.
\end{itemize}

Szczegółowe kryteria zakończenia epizodu zostaną doprecyzowane po analizie możliwości wykrywania poszczególnych zdarzeń za pomocą telemetrii.

\subsection{Założenia i ograniczenia środowiska}

W projekcie przyjmuje się następujące wstępne założenia:
\begin{itemize}
    \item agent będzie sterował pojazdem w grze TrackMania 2020,
    \item komunikacja z grą będzie realizowana z wykorzystaniem biblioteki TrackmaniaRL,
    \item agent będzie korzystał z danych telemetrycznych zamiast bezpośredniej analizy obrazu z ekranu,
    \item uczenie będzie odbywać się poprzez wielokrotne przejazdy wyznaczonej trasy,
    \item podstawowym celem będzie uzyskanie możliwie krótkiego czasu przejazdu przy zachowaniu skuteczności pokonywania trasy.
\end{itemize}

Do najważniejszych ograniczeń należą fizyka pojazdu, geometria trasy, ograniczenia sterowania oraz dostępność i częstotliwość danych telemetrycznych.

\subsection{Pojedynczy krok interakcji}

Pojedynczy krok interakcji agenta ze środowiskiem będzie obejmował:
\begin{enumerate}
    \item pobranie aktualnej obserwacji telemetrycznej,
    \item wybór akcji przez agenta,
    \item przekazanie akcji do środowiska gry,
    \item wykonanie ruchu pojazdu,
    \item pobranie nowej obserwacji,
    \item obliczenie nagrody za wykonane działanie,
    \item sprawdzenie warunków zakończenia epizodu.
\end{enumerate}

W ten sposób agent będzie stopniowo zdobywał doświadczenie i doskonalił strategię sterowania.

\subsection{Kryteria oceny skuteczności}

Skuteczność agenta będzie oceniana na podstawie następujących miar:
\begin{itemize}
    \item średni czas ukończenia trasy,
    \item odsetek epizodów zakończonych ukończeniem trasy,
    \item średni postęp uzyskiwany podczas epizodu,
    \item liczba kolizji i przypadków wypadnięcia poza trasę,
    \item stabilność uzyskiwanych wyników pomiędzy kolejnymi przejazdami.
\end{itemize}

Wyniki agenta uczącego się będą mogły zostać porównane z wynikami agenta bazowego, aby ocenić rzeczywistą skuteczność zastosowanego algorytmu.

\subsection{Elementy wymagające weryfikacji}

Po pierwszych testach środowiska mogą wymagać korekty następujące elementy formalizacji:
\begin{itemize}
    \item dostępność i dokładna struktura danych telemetrycznych,
    \item wybór zmiennych tworzących obserwację agenta,
    \item sposób reprezentacji przestrzeni akcji,
    \item wartości współczynników funkcji nagrody,
    \item kryteria wykrywania kolizji i wypadnięcia poza trasę,
    \item warunki zakończenia epizodu,
    \item częstotliwość pobierania obserwacji i wykonywania akcji.
\end{itemize}

Ostateczna formalizacja zostanie dostosowana do rzeczywistych możliwości środowiska oraz wyników pierwszych testów.

\end{document}
```